\documentclass[10pt, a4paper]{article}
% ===== Essential Packages =====
\usepackage[utf8]{inputenc}       % UTF-8 encoding
\usepackage[T1]{fontenc}          % Better font encoding
\usepackage{lmodern}              % Modern Latin Modern font
\usepackage{ctex}                 % Chinese support (must be loaded after fontenc)
\usepackage{subcaption}
% ===== Math & Symbols =====
\usepackage{amsmath}              % Advanced math
\usepackage{amssymb}              % Additional math symbols
\usepackage{mathrsfs}             % Ralph Smith's Formal Script
\usepackage{braket}               % Dirac bra-ket notation

% ===== Graphics & Tables =====
\usepackage{graphicx}             % Image support
\usepackage{float}                % Better float control
\usepackage{wrapfig}              % Wrapped figures
\usepackage{tikz}                 % Vector graphics
\usepackage[table]{xcolor}        % Colored tables
\usepackage{array}                % Enhanced tables
\usepackage{multirow}             % Multirow tables
\usepackage{tabularx}             % Adjustable-width tables
\usepackage{booktabs}             % Professional table lines
\usepackage{siunitx}              % SI units in tables
\usepackage{caption}              % Better captions
\usepackage{adjustbox}            % Box resizing

% ===== Code Listings =====
\usepackage{listings}             % Code blocks
\usepackage{tcolorbox}            % Colored boxes
\tcbuselibrary{listings, breakable} % Listings in tcolorbox

% ===== Misc =====
\usepackage{lipsum}               % Dummy text
\usepackage{footmisc}             % Better footnotes
\usepackage[margin=1in]{geometry} % Page margins
\usepackage{diagbox}
\usepackage{ctex}
\author{张世航}
\title{分光计测量物质折射率}
\begin{document}
	\maketitle
	\section{测量三棱镜的折射率}
	\begin{tabular}{|c|c|c|c|c|}
		\hline
		$\varphi_1'$&$\varphi_2'$  & $\varphi_1''$ & $\varphi_2''$&$\delta_{min}^1$ \\
		\hline
		$248^o37'$&$209^o56'$  & $68^o37'$ & $29^o56'$ &$ 38^o41'$\\
		\hline
	\end{tabular}\\
	
	\begin{tabular}{|c|c|c|c|c|}
		\hline
		$\varphi_1'$&$\varphi_2'$  & $\varphi_1''$ & $\varphi_2''$ &$\delta_{min}^2$\\
		\hline
		$198^o54'$&$160^o10'$  & $18^o54'$ & $340^o10'$ &$38^o44'$\\
		\hline
	\end{tabular}\\
	
	\begin{tabular}{|c|c|c|c|c|}
		\hline
		$\varphi_1'$&$\varphi_2'$  & $\varphi_1''$ & $\varphi_2''$ &$\delta_{min}^3$\\
		\hline
		$216^o33'$&$177^o48'$  & $36^o33'$ & $357^o48'$ &$38^o45'$\\
		\hline
	\end{tabular}
	
	\[\delta_{min} = \frac{(\varphi_2'-\varphi_1')+(\varphi_2''-\varphi_1'')}{2}\]
	\[\bar{\delta}_{min} = \frac{\delta_{min}^1+\delta_{min}^2+\delta_{min}^3}{3}\approx38^o43'\]
	\[n_{\text{三棱镜}} = \frac{\sin\frac{\bar{\delta}_{min}+A}{2}}{\sin\frac{A}{2}} =1.5176588530685438\]
	由于我们三次测量相互独立(即我们测量一次就转了一次刻度盘),我们没办法使用间接测量的误差传递公式进行计算.\\
	我们只能对三次计算出的$\delta_{min}$进行A类不确定度分析.\\
	\[u_A(\delta_{min}) = \sqrt{\frac{\sum_{i}^{N}(\delta_{min}^i-\bar{\delta}_{min})^2}{N(N-1)}}=0.020030840419245546^o\]
	\[u_B(\delta_{min}) = \frac{\Delta}{\sqrt{3}} = \frac{\frac{1}{60}^o}{\sqrt{3}} = 0.009622504486493762
	^o\]
	\[u_{\delta_{min}} = np.sqrt(u_A^2(\delta_{min})+u_B^2(\delta_{min})) = 0.022222222222223267^o\]
	\[U = \frac{\partial n}{\partial \delta_{min}}u_{\delta_{min}} = \frac{1}{\sin\frac{A}{2}}\cos\frac{\delta_{min}+A}{2}\times\frac{1}{2}u_{\delta_{min}} =0.000387850944887647\]
	则\[n = 1.5177\pm0.0004\]
	\section{测量水的折射率}
	\begin{tabular}{|c|c|c|c|c|}
		\hline
		$\text{分界线}_1$&$\text{十字叉丝}_1$  & $\text{分界线}_2$&$\text{十字叉丝}_2$&$r'$\\
		\hline
		$154^o32'$&$156^o47'$  & $334^o32'$ & $336^o47'$  &$ 2^o15' $\\
		\hline
	\end{tabular}\\
	
	\begin{tabular}{|c|c|c|c|c|}
		\hline
		$\text{分界线}_1$&$\text{十字叉丝}_1$  & $\text{分界线}_2$&$\text{十字叉丝}_2$&$r'$\\
		\hline
		$175^o47'$&$178^o2'$  & $355^o47'$ & $358^o2'$ &$ 2^o15' $\\
		\hline
	\end{tabular}\\
	
	\begin{tabular}{|c|c|c|c|c|}
		\hline
		$\text{分界线}_1$&$\text{十字叉丝}_1$  & $\text{分界线}_2$&$\text{十字叉丝}_2$&$r'$\\
		\hline
		$197^o47'$&$200^o2'$  & $17^o47'$ & $20^o2'$ &$ 2^o15' $\\
		\hline
	\end{tabular}\\
	\[\bar{r}' = 2^o15'\]
	\[\bar{n}_x = \sin{A}\sqrt{n^2 - \sin^2r'} + \cos{A}sinr' = 1.3335211879605458\]
	对于$r'$,由于三次测量的结果相同,则A类不确定度不存在.\\
	\[u_{r'} = u_B(r') =  \frac{\Delta}{\sqrt{3}} = \frac{\frac{1}{60}^o}{\sqrt{3}} = 0.009622504486493762
	^o\]
	\[U = \sqrt{(\frac{\partial n_x}{\partial \delta_{min}})^2\times u^2(\delta_{min})+(\frac{\partial n_x}{\partial r'})^2\times u^2(r')}\]
	\[\frac{\partial n_x}{\partial \delta_{min}}u_{\delta_{min}} = \sin{A}\frac{n}{\sqrt{n^2-\sin^2r'}}\frac{1}{\sin\frac{A}{2}}\cos\frac{\delta_{min}+A}{2}\times\frac{1}{2}u_{\delta_{min}} = 0.0002188340339156575\]
	\[\frac{\partial n_x}{\partial r'}u_{r'} =  [sin{A}\frac{-\sin{r'}\cos{r'}}{\sqrt{n^2-\sin^2{r'}}}+\cos{A}\cos{r'}]u_{r'} = 7.819983676468333\times10^{-5}\]
	\[u_{n_x} = \sqrt{0.01253826655718989^2+0.004480520605228325^2} = 0.00023238663659905717\]
	\[n_x = 1.3335\pm 0.0002\]
\end{document}